% Historical note on the abandoned reflection route.

This chapter records a second possible route to the plus-to-minus divisibility
step.  It is not the route used in the proof of Kummer's criterion in
\cref{ch:kummer-final}.  Its purpose here is to isolate the remaining
mathematical input and to explain how the reflection argument would proceed if
that input were supplied.

\section{The aim}

Let \(p\) be an odd prime and let \(K=\Q(\zeta_p)\).  Put
\[
  A=\Cl{\mathcal O_K}/p\Cl{\mathcal O_K},
\]
with the natural action of
\(\Delta=\Gal(K/\Q)\simeq(\Z/p\Z)^\times\).  The reflection route tries to
prove
\[
  p\mid h^+(K)\quad\Longrightarrow\quad p\mid h^-(K)
\]
by proving a statement about the character components of \(A\).  If
\(A_i\) denotes the \(i\)-th eigenspace, the desired reflection statement is
that a nonzero even component forces a nonzero reflected component:
\[
  A_i\ne0 \quad\Longrightarrow\quad A_{1-i}\ne0.
\]
The index \(1-i\) is always understood modulo \(p-1\), in the standard
representative range.

\section{The unproved reciprocity input}

\begin{theorem}[One-sided Kummer principal reciprocity]
  \label{thm:abandoned-oskr}
  \lean{BernoulliRegular.Furtwaengler.oneSidedKummerPrincipalReciprocity_canonical}
  \notready
Let \(\alpha,\beta\in\mathcal O_K\) be nonzero, coprime, and prime to \(p\).
Assume that \(\alpha\) is a \(p\)-th power in the
\((\zeta_p-1)\)-adic completion of \(K\).  Then the canonical
\(p\)-th-power residue symbols satisfy
\[
  \left(\frac{\alpha}{(\beta)}\right)_p
  =
  \left(\frac{\beta}{(\alpha)}\right)_p.
\]
\end{theorem}

This is the single missing mathematical input for the reflection route as
recorded here.  It is intentionally stated as a concrete reciprocity theorem,
not as a packaged reflection principle.

The immediate consequence needed below is the following principal-denominator
vanishing statement.  If \(\eta\) is locally a \(p\)-th power at
\(\zeta_p-1\), prime to \(p\), and
\[
  (\eta)=\mathfrak b^p,
\]
then for every admissible principal denominator \((\gamma)\) one has
\[
  \left(\frac{\eta}{(\gamma)}\right)_p=1.
\]
Indeed, reciprocity gives
\[
  \left(\frac{\eta}{(\gamma)}\right)_p
  =
  \left(\frac{\gamma}{(\eta)}\right)_p,
\]
and the denominator \((\eta)=\mathfrak b^p\) makes the right hand side
trivial by multiplicativity in the denominator.

\section{Component reflection}

\begin{theorem}[Component reflection]
  \label{thm:abandoned-component-reflection}
  \lean{BernoulliRegular.weakReflection_componentNontrivial}
  \leanok
  \uses{thm:abandoned-oskr}
Let \(i\) be an even reflection index.  If the \(i\)-component of
\(\Cl{\mathcal O_K}/p\Cl{\mathcal O_K}\) is nonzero, then the reflected
component is nonzero:
\[
  A_i\ne0\quad\Longrightarrow\quad A_{1-i}\ne0.
\]
\end{theorem}

\begin{proof}\leanok
The proof has five mathematical steps.

First, the nonzero component \(A_i\) gives a nonzero locally-primary singular
pseudo-unit.  More precisely, one obtains
\[
  \eta\in\mathcal O_K,\qquad
  (\eta)=\mathfrak b^p,\qquad
  \eta\notin K^{\times p},
\]
where \(\eta\) is prime to \(p\), locally a \(p\)-th power at
\(\zeta_p-1\), and has Galois weight \(i\) up to multiplication by global
\(p\)-th powers.

Second, use the residue symbol to define a character on ideal classes:
\[
  \chi_\eta([\mathfrak a])
  =
  \left(\frac{\eta}{\mathfrak a}\right)_p.
\]
The representative \(\mathfrak a\) is chosen coprime to a finite bad set
containing the primes above \(p\), the primes dividing \((\eta)\), and the
Kummer-Dedekind conductor.  Such representatives exist by ideal avoidance.
The principal-denominator vanishing consequence of
\cref{thm:abandoned-oskr} makes the value independent of this choice.  Thus
\(\chi_\eta\) descends to a character of \(A\).

Third, the character is nonzero.  If it were zero, then for every prime ideal
\(\mathfrak q\) outside the bad set the residue symbol
\((\eta/\mathfrak q)_p\) would be trivial.  Hence the reduction of \(\eta\)
would be a \(p\)-th power in \(\mathcal O_K/\mathfrak q\), so
\(T^p-\eta\) would split modulo \(\mathfrak q\).  Kummer-Dedekind would then
show that almost all primes of \(K\) split completely in
\(K(\eta^{1/p})/K\).  The weak splitting lemma forces this extension to be
trivial, contradicting \(\eta\notin K^{\times p}\).

Fourth, the Galois weight of \(\eta\) gives the covariance relation
\[
  \chi_\eta(\sigma_a x)=a^{1-i}\chi_\eta(x)
  \qquad (a\in(\Z/p\Z)^\times).
\]
Thus the nonzero character is supported on the reflected component.

Finally, the boundary case \(i=p-1\) is treated separately.  The zero-character
component of \(A\) is trivial: if an ideal class is fixed by all Galois
conjugates, the product of all its conjugates is represented by a
Galois-stable ideal, hence by a principal ideal.  The class therefore has
order dividing both \(p\) and \(p-1\), so it is trivial.  The remaining case
is the interior even case handled by the preceding argument.
\end{proof}

\section{From reflection to class numbers}

\begin{theorem}[Reflection-route plus-to-minus implication]
  \label{thm:abandoned-reflection-plus-minus}
  \lean{BernoulliRegular.weakReflection_dvd_hMinus_of_dvd_hPlus}
  \leanok
  \uses{thm:abandoned-component-reflection}
For an odd prime \(p\),
\[
  p\mid h^+(K)\quad\Longrightarrow\quad p\mid h^-(K).
\]
\end{theorem}

\begin{proof}\leanok
Assume \(p\mid h^+(K)\) and \(p\nmid h^-(K)\).  Since
\[
  h(K)=h^+(K)h^-(K),
\]
the first assumption implies \(p\mid h(K)\), so the group \(A\) is nonzero.
Decompose \(A\) into character components and choose a nonzero component
\(A_k\).

The assumption \(p\nmid h^-(K)\) says that complex conjugation acts trivially
on the relevant mod-\(p\) class group.  Therefore no odd component can be
nonzero.  The zero component is also trivial, as in the boundary argument
above.  Hence the chosen nonzero component is even and nonzero.  Component
reflection gives a nonzero reflected component, but the reflected index is
odd, contradicting the trivial action of complex conjugation.  Thus
\(p\nmid h^-(K)\) is impossible.
\end{proof}

Once the plus-to-minus implication is known, the reflection route also
recovers the same class-number criterion as the main proof:
\[
  (p:\N)\mid h(K)
  \quad\Longleftrightarrow\quad
  \exists k,\ 1\le k,\ 2k\le p-3,\quad
  (p:\Z)\mid (B_{2k})_{\mathrm{num}}.
\]

Indeed, by the factorisation \(h(K)=h^+(K)h^-(K)\), divisibility of \(h(K)\)
by the
prime \(p\) means that \(p\) divides either \(h^+(K)\) or \(h^-(K)\).  The
plus-to-minus implication transfers the first case to the second.  Thus
\[
  p\mid h(K)\quad\Longleftrightarrow\quad p\mid h^-(K).
\]
The minus class-number criterion from \cref{thm:hminus-bernoulli-final}
identifies the latter condition with Bernoulli numerator divisibility in
Kummer's range.

\section{How the reciprocity input is to be proved}

The intended source for \cref{thm:abandoned-oskr} is the classical
Furtwaengler--Kummer reciprocity argument, in the form developed by
Ireland--Rosen.  The already completed part of that source proves the
integer-denominator Eisenstein reciprocity theorem:

\begin{theorem}[Ireland--Rosen integer reciprocity]
  \label{thm:ireland-rosen-integer-reciprocity}
  \lean{BernoulliRegular.Furtwaengler.IrelandRosen.eisensteinReciprocity_theorem1}
  \leanok
Let \(\alpha\in\mathcal O_K\) be primary and prime to \(p\), and let
\(a\in\Z\) be nonzero, prime to \(p\), and coprime to \((\alpha)\).  Then
\[
  \left(\frac{\alpha}{(a)}\right)_p
  =
  \left(\frac{a}{(\alpha)}\right)_p .
\]
\end{theorem}

\begin{proof}\leanok
The proof follows the Gauss--Jacobi sum construction.  For every prime ideal
\(P\nmid p\), one constructs an element \(\Phi(P)\) whose principal ideal is
controlled by the Stickelberger element:
\[
  (\Phi(P))=P^\theta .
\]
This is proved without assuming that \(P\) has residue degree one.  The
construction also supplies the semiprimary normalization, the conjugation-norm
identity, and the Frobenius residue-symbol identity for \(\Phi(P)\).

For a general principal ideal \((\alpha)\), take the product over the prime
factorization of \((\alpha)\):
\[
  \Phi(\alpha)=
  \prod_{P\mid(\alpha)}\Phi(P)^{v_P((\alpha))}.
\]
The Stickelberger identities multiply to give
\[
  (\Phi(\alpha))=(\alpha^\theta).
\]
Hence \(\Phi(\alpha)=\varepsilon\,\alpha^\theta\) for a unit
\(\varepsilon\).  The primary-unit lemma says that when \(\alpha\) is
primary, this unit contributes only a sign, and a sign has trivial
\(p\)-th-power residue symbol because \(p\) is odd.

Combining the product identity, the primary-unit cancellation, and the
denominator norm calculation gives the ideal-norm reciprocity formula.  Taking
the denominator ideal to be \((a)\) gives the displayed integer reciprocity
law.
\end{proof}

To finish the reflection route one still has to pass from this
integer-denominator theorem to the full one-sided principal reciprocity
statement of \cref{thm:abandoned-oskr}.  The missing step is not a new
reflection theorem; it is the remaining principal-denominator reciprocity
argument that allows a general coprime principal denominator \((\beta)\) in
place of a rational integer denominator \((a)\), with the local-primary
hypothesis on the numerator preserved throughout.
